%%%%%%%%%%%%%%%%%%%%%%%%%%%%%%%%%%%%%%%%%%%%%%%%%%%%%%%%%%%%%
%% Begin exercise %%
%%%%%%%%%%%%%%%%%%%%%%%%%%%%%%%%%%%%%%%%%%%%%%%%%%%%%%%%%%%%%
\ex{Induction machines}


\normalsize{\textbf{Acknowledgement}: The following exercise is adapted from ``Geregelte Drehstromantriebe / Controlled AC Drives'' by J. Böcker, Paderborn University, 2021
}\\


%%%%%%%%%%%%%%%%%%%%%%%%%%%%%%%%%%%%%%%%%%%%%%%%%%%%%%%%%%%%%
%% Task 1 %%
%%%%%%%%%%%%%%%%%%%%%%%%%%%%%%%%%%%%%%%%%%%%%%%%%%%%%%%%%%%%%

\task{Steady-state operation of an induction machine}

An induction machine with the characteristics in Tab.~\ref{tab:characteristicsIM_task2} is given.
\begin{table}[htb]
    \caption{Characteristics of the given induction machine.}
    \centering
    \begin{tabular}{lll}\toprule
    Symbol  & Description       & Values \\
    \midrule
    $U_{\mathrm{n}}$    & Nominal voltage           & $\SI{380}{\volt}$ \\
    $I_{\mathrm{n}}$    & Nominal phase current     & $\SI{54}{\ampere}$ \\
    $f_{\mathrm{s,n}}$  & Nominal frequency         & $\SI{50}{\hertz}$ \\
    $P_{\mathrm{n}}$    & Nominal power             & $\SI{25}{\kilo\watt}$ \\
    $n_{\mathrm{n}}$    & Nominal speed             & $\SI{1465}{\per\minute}$ \\
    $\cos(\varphi)$     & Power factor              & 0.77 \\
    \midrule
    $R_{\mathrm{s}}$    & Stator resistance         & $\SI{0.48}{\Omega}$ \\
    $R_{\mathrm{r}}'$    & Rotor resistance          & $\SI{85}{\milli\Omega}$ \\
    $M$                 & Mutual inductance         & $\SI{100}{\milli\henry}$ \\
    $L_{\mathrm{\sigma,s}}$    & Stator leakage inductance  & $\SI{2}{\milli\henry}$ \\
    $L_{\mathrm{\sigma,r}}'$    & Rotor leakage inductance   & $\SI{2}{\milli\henry}$ \\
    \bottomrule
    \end{tabular}
    \label{tab:characteristicsIM_task2}
\end{table}


%%%%%%%%%%%%%%%%%%%%%%%%%%%%%%%%%%%%%%%%%%%%%%%%%%%%%%%%%%%%%
\subtask{Determine the amplitude of the complex AC voltage $\hat{U}$ and current phasor $\hat{I}$.}

\begin{solutionblock}
    The amplitude of the line-to-line voltage is given with
    \begin{equation}
    \hat{U} = \sqrt{2}\cdot \SI{380}{\volt}
    = \SI{537.4}{\volt},
    \end{equation}
    and, therefore, the amplitude of the phase-to-star-point voltage is defined as:
    \begin{equation}
    \hat{U}_{\mathrm{star}} = \frac{\SI{380}{\volt}}{\sqrt{3}}\cdot \sqrt{2} = \SI{310.9}{\volt}.
    \end{equation}

    Furthermore, the current amplitude is defined by:
    \begin{equation}
        \hat{I} = \sqrt{2} I = \sqrt{2} \cdot \SI{54}{\ampere} = \SI{76.4}{\ampere}.
    \end{equation}
\end{solutionblock}

%%%%%%%%%%%%%%%%%%%%%%%%%%%%%%%%%%%%%%%%%%%%%%%%%%%%%%%%%%%%%
\subtask{Determine the number of pole pairs $p$.}

\begin{solutionblock}
    The stator frequency is given with $f_{\mathrm{s}}$ = 50 Hz, which leads to the synchronous rotational speed of the electrical frequency as follows:
    \begin{equation}
        n_{\mathrm{syn,el}} = \SI{50}{\hertz} \cdot \SI{60}{\second\per\minute}
        = \SI{3000}{\per\minute}.
    \end{equation}
    
    The nominal speed is given with $n_{\mathrm{n}} = \SI{1465}{\per\minute}$, which is associated with a synchronous speed of $n_{\mathrm{syn,mech}} = \SI{1500}{\per\minute}$.
    
    Therefore, the number of pole pairs is defined with:
    \begin{equation}
        p = \frac{n_{\mathrm{syn,el}}}{n_{\mathrm{syn,mech}}}
        = \frac{\SI{3000}{\per\minute}}{\SI{1500}{\per\minute}}
        = 2.
    \end{equation}

\end{solutionblock}

%%%%%%%%%%%%%%%%%%%%%%%%%%%%%%%%%%%%%%%%%%%%%%%%%%%%%%%%%%%%%
\subtask{Calculate the rated slip $s_{\mathrm{n}}$. Which slip occurs at an ideal no-load operation (no firction)?}

\begin{solutionblock}
    The slip is calculated by:
    \begin{equation}
        s_{\mathrm{n}} = \frac{\omega_{\mathrm{slip,n}}}{\omega_{\mathrm{s}}}
        = \frac{\omega_{\mathrm{s}}-p \omega_{\mathrm{r}}}{\omega_{\mathrm{s}}}
        = \frac{\SI{2\pi\cdot \frac{3000}{60}}{\per\second}- 2 \cdot \SI{2\pi \cdot \frac{1465}{60}}{\per\second}}{\SI{2\pi \cdot \frac{3000}{60}}{\per\second}}
        = 0.023,
    \end{equation}
    where $\omega_{\mathrm{r}}$ is the rotor speed.
    At an ideal no-load operation (also no friction), the rotor has the same speed as the stator field and, therefore, $s$ = 0.
\end{solutionblock}

%%%%%%%%%%%%%%%%%%%%%%%%%%%%%%%%%%%%%%%%%%%%%%%%%%%%%%%%%%%%%
\subtask{Calculate the apparent, the reactive and the electrical power. In addition, determine the efficiency $\eta_{\mathrm{n}}$ for the rated operating point.}

\begin{solutionblock}
    The apparent power is calculated with
    \begin{equation}
        S = \sqrt{3} U I
        = \sqrt{3} \cdot \SI{380}{\volt} \cdot \SI{54}{\ampere}
        = \SI{35.541}{\kilo\volt\ampere},
    \end{equation}
    
    and, the electrical power is determined with the power factor in the task as follows:
    \begin{equation}
        P_{\mathrm{el,n}} = \sqrt{3} U I \cos(\varphi)
        = \sqrt{3} \cdot \SI{380}{\volt} \cdot \SI{54}{\ampere} \cdot 0.77
        = \SI{27.367}{\kilo\watt}.
    \end{equation}

    The angle of the phase shift is calculated by:
    \begin{equation}
        \varphi = \cos^{-1}(0.77) = \SI{39.6}{\degree}.
    \end{equation}

    Hence, the reactive power yields:
    \begin{equation}
        Q = \sqrt{3} U I \sin(\varphi)
        = \sqrt{3} \cdot \SI{380}{\volt} \cdot \SI{54}{\ampere} \cdot \sin(\SI{39.6}{\degree})
        = \SI{22.655}{\kilo\volt\ampere}.
    \end{equation}


    The efficiency results as:
    \begin{equation}
        \eta_{\mathrm{n}} = \frac{P_{\mathrm{mech,n}}}{P_{\mathrm{el,n}}} = \frac{\SI{25}{\kilo\watt}}{\SI{27.367}{\kilo\watt}}
        = 0.914.
    \end{equation}
\end{solutionblock}

%%%%%%%%%%%%%%%%%%%%%%%%%%%%%%%%%%%%%%%%%%%%%%%%%%%%%%%%%%%%%
\subtask{Determine the nominal torque generated by the induction machine.}

\begin{solutionblock}
    The torque is calculated as:
    \begin{equation}
        T_{\mathrm{n}} = \frac{P_{\mathrm{mech}}}{\omega_{\mathrm{mech}}}
        = \frac{\SI{25}{\kilo\watt}}{\SI{2\pi\cdot \frac{1465}{60}}{\per \second}}
        = \SI{163}{\newton \metre}.
    \end{equation}
\end{solutionblock}

%%%%%%%%%%%%%%%%%%%%%%%%%%%%%%%%%%%%%%%%%%%%%%%%%%%%%%%%%%%%%
\subtask{Calculate the starting $T_{\mathrm{0}}$ and the maximum torque $T_{\mathrm{max}}$. For the latter determine first the slip $s_{\mathrm{max}}$ at the operating point with the maximum torque.}

\begin{solutionblock}
    The slip at the maximum torque operating point is defined by
    \begin{equation}
        s_{\mathrm{max}} = \frac{R_{\mathrm{r}}'}{\sigma \left(L_{\mathrm{\sigma,r}}' + M\right) \omega_{\mathrm{s}}}
        = \frac{\SI{85}{\milli\Omega}}{0.038 \cdot (\SI{2}{\milli\henry}+\SI{100}{\milli\henry}) \cdot \SI{2\pi\cdot 50}{\per\second}}
        = 0.068,
    \end{equation}
    with the leakage coefficient:
    \begin{align}
        \begin{split}
            \sigma &= 1- \frac{M^2}{(M+L_{\mathrm{\sigma,s}})(M+L_{\mathrm{\sigma,r}}')}\\
            &= 1 - \frac{\left(\SI{100}{\milli\henry}\right)^2}{\left(\SI{102}{\milli\henry}\right)\cdot\left(\SI{102}{\milli\henry}\right)}\\
            &= 0.038.
        \end{split}
    \end{align}

    Hence, the maximum torque is calculated with:
    \begin{align}
        \begin{split}
            T_{\mathrm{max}} &= \frac{3}{2}p \frac{U_{\mathrm{s}}^2}{\omega_{\mathrm{s}}^2}\frac{M^2}{\sigma \left(L_{\mathrm{\sigma,s}}+M \right)^2\left(L_{\mathrm{\sigma,r}}'+M\right)} \\
            &= \frac{3}{2}\cdot2\cdot\frac{\frac{\SI{380}{\volt}}{\sqrt{3}}}{\left(\SI{2\pi\cdot 50}{\per\second}\right)^2} \cdot \frac{\SI{100}{\milli\henry}}{0.038\cdot\left(\SI{102}{\milli\henry}\right)^2 \cdot \left(\SI{102}{\milli\henry}\right)}\\
            &= \SI{355}{\newton\metre}.
        \end{split}
    \end{align}

    The starting torque is determined as follows:
    \begin{equation}
        T_{\mathrm{0}} = T_{\mathrm{max}} \frac{2 s_{\mathrm{max}}}{1 + s_{\mathrm{max}}^2}
        = \SI{355}{\newton\metre}\cdot \frac{2 \cdot 0.068}{1 + 0.068^2} = \SI{48.3}{\newton\metre}.
    \end{equation}
\end{solutionblock}


%%%%%%%%%%%%%%%%%%%%%%%%%%%%%%%%%%%%%%%%%%%%%%%%%%%%%%%%%%%%%
\subtask{Determine the stator currents $I_{\mathrm{s,dq}}$, and, in addition, the induced rotor currents $I_{\mathrm{r,dq}}$.}

\begin{solutionblock}
    Derived from the lecture notes, the stator currents are calculated as follows:
    \begin{align}
        \begin{split}
        I_{\mathrm{s,d}} &= \frac{U_{\mathrm{s}}}{\omega_{\mathrm{s}}} \frac{\sigma^2\omega_{\mathrm{slip}}^2 \left(L_{\mathrm{\sigma,s}}+M\right)\left(L_{\mathrm{\sigma,r}}'+M\right)^3 + \left(L_{\mathrm{\sigma,r}}'+M\right)\left(L_{\mathrm{\sigma,s}}+M\right)\left(R_{\mathrm{r}}'\right)^2 - M^2 \left(R_{\mathrm{r}}'\right)^2}{\sigma \left(L_{\mathrm{\sigma,s}}+M\right)^2 \left(L_{\mathrm{\sigma,r}}'+M\right)\omega_{\mathrm{slip}}\left(\sigma^2 \omega_{\mathrm{slip}}^2 \left(L_{\mathrm{\sigma,r}}'+M\right)^2 + \left(R_{\mathrm{r}}'\right)^2\right)}\\
        &= \SI{3.3}{\ampere},
        \end{split}
    \end{align}
    
    \begin{equation}
        I_{\mathrm{s,q}} = \frac{U_{\mathrm{s}}}{\omega_{\mathrm{s}}} \frac{M^2}{\sigma\left(L_{\mathrm{\sigma,s}}+M\right)^2 \left(L_{\mathrm{\sigma,r}}'+M\right)} \frac{1}{\frac{\omega_{\mathrm{slip}}}{\omega_{\mathrm{max}}} + \frac{\omega_{\mathrm{max}}}{\omega_{\mathrm{slip}}}}
        = \SI{51.8}{\ampere}.
    \end{equation}
    
    
    The rotor currents are calculated as follows:
    \begin{equation}
        I_{\mathrm{r,d}} = -U_{\mathrm{s}} \frac{M s}{\left(L_{\mathrm{\sigma,s}}+M\right) R_{\mathrm{r}}'}\frac{1}{\frac{\omega_{\mathrm{slip}}}{\omega_{\mathrm{max}}} + \frac{\omega_{\mathrm{max}}}{\omega_{\mathrm{slip}}}}
        = \SI{-18.06}{\ampere},
    \end{equation}

    \begin{equation}
        I_{\mathrm{r,q}} = -\frac{U_{\mathrm{s}}}{\omega_{\mathrm{s}}}\frac{M}{\sigma \left(L_{\mathrm{\sigma,s}}+M\right)\left(L_{\mathrm{\sigma,r}}'+M\right)}\frac{1}{\frac{\omega_{\mathrm{slip}}}{\omega_{\mathrm{max}}} + \frac{\omega_{\mathrm{max}}}{\omega_{\mathrm{slip}}}}
        = \SI{-52.9}{\ampere}.
    \end{equation}

    

\end{solutionblock}


%%%%%%%%%%%%%%%%%%%%%%%%%%%%%%%%%%%%%%%%%%%%%%%%%%%%%%%%%%%%%
\subtask{Calculate the losses in the stator winding and in the rotor.}

\begin{solutionblock}
    The stator current is determined by
    \begin{equation}
        I_{\mathrm{s,dq}} = \sqrt{I_{\mathrm{s,d}}^2 + I_{\mathrm{s,q}}^2}
        = \sqrt{\left(\SI{3.3}{\ampere}\right)^2 + \left(\SI{51.8}{\ampere}\right)^2}
        = \SI{51.9}{\ampere},
    \end{equation}
    which results to the stator winding losses as follows:
    \begin{equation}
        P_{\mathrm{l,stator}} = \frac{3}{2} I_{\mathrm{s,dq}}^2 R_{\mathrm{s}}
        = \frac{3}{2} \cdot \left(\SI{51.9}{\ampere} \right)^2 \cdot \SI{0.48}{\Omega}
        = \SI{1939.4}{\watt}.
    \end{equation}

    The rotor current is calculated in the same way. This leads to
    \begin{equation}
        I_{\mathrm{r,dq}} = \sqrt{I_{\mathrm{r,d}}^2 + I_{\mathrm{r,q}}^2}
        = \sqrt{\left(\SI{-18.1}{\ampere}\right)^2 + \left(\SI{-52.9}{\ampere}\right)^2}
        = \SI{55.9}{\ampere},
    \end{equation}
    and the ohmic rotor losses are calculated with:
    \begin{equation}
        P_{\mathrm{l,rotor}} = \frac{3}{2} I_{\mathrm{r,dq}}^2 R_{\mathrm{r}}'
        = \frac{3}{2} \cdot \left(\SI{55.9}{\ampere} \right)^2 \cdot \SI{85}{\milli\Omega}
        = \SI{398.4}{\watt}.
    \end{equation} 


\end{solutionblock}


%%%%%%%%%%%%%%%%%%%%%%%%%%%%%%%%%%%%%%%%%%%%%%%%%%%%%%%%%%%%%
\subtask{Draw the space vector diagram for the operation under rated conditions including the vectors $\underline{u}_{\mathrm{s}}$, $\underline{i}_{\mathrm{s}}$, $\underline{\psi}_{\mathrm{s}}$ and $\frac{\mathrm{d}}{\mathrm{d}t} \underline{\psi}_{\mathrm{s}}$. Assume that the current vector $\underline{i}_{\mathrm{s}}$ is oriented along the x-axis of the Cartesian coordinate system.}

\begin{solutionblock}
    The current is aligned with the x-axis of the Cartesian coordinate system. Therefore, the current is given with:
    \begin{equation}
        \underline{i}_{\mathrm{s}} = \sqrt{2} I e^{\mathrm{j}0^{\circ}}
        = \SI{76.37}{\ampere} \cdot e^{\mathrm{j}0^{\circ}}.
    \end{equation}

    The angle shift between the voltage and current is calculated by
    \begin{equation}
        \varphi_{\mathrm{ui}} = \SI{39.65}{\degree},
    \end{equation}

    therefore, the voltage is defined as follows:
    \begin{equation}
        \underline{u}_{\mathrm{s}} = \sqrt{2} U e^{\mathrm{j} \varphi_{\mathrm{ui}}}
        = \SI{310.27}{\volt} \cdot e^{\mathrm{j}39.65^{\circ}}
        = \SI{238.89}{\volt} + \mathrm{j} \SI{197.98}{\volt}.
    \end{equation}

    The stator resistance $R_{\mathrm{s}}$ is given in the task. Hence, the voltage equation is defined as:
    \begin{equation}
        \underline{u}_{\mathrm{s}} = R_{\mathrm{s}} \underline{i}_{\mathrm{s}} + \frac{\mathrm{d}}{\mathrm{d}t}\psi_{\mathrm{s}}.
    \end{equation}

    Rearranging leads to:
    \begin{align}
        \begin{split}
        \frac{\mathrm{d}}{\mathrm{d}t}\psi_{\mathrm{s}} &= \underline{u}_{\mathrm{s}} - R_{\mathrm{s}} \underline{i}_{\mathrm{s}} \\
        &= \SI{238.89}{\volt} + \mathrm{j} \SI{197.98}{\volt} - \SI{0.48}{\Omega} \cdot \SI{76.37}{\ampere} \\
        &= \SI{202.23}{\volt} + \mathrm{j} \SI{197.98}{\volt} \\
        &= \SI{283}{\volt} \cdot e^{\mathrm{j}44.4^{\circ}}.
        \end{split}
    \end{align}

    The stator flux vector is defined as
    \begin{equation}
        \underline{\psi}_{\mathrm{s}}(t) = \psi_{\mathrm{s}}(t) \cdot e^{\mathrm{j}\varphi_{\mathrm{s}}(t)},
    \end{equation}

    and the differentiation (under consideration of the product rule) results in:
    \begin{equation}
        \frac{\mathrm{d}}{\mathrm{d}t}\underline{\psi}_{\mathrm{s}}(t)
        = \frac{\mathrm{d}}{\mathrm{d}t} \psi_{\mathrm{s}}(t) \cdot e^{\mathrm{j}\varphi_{\mathrm{s}}(t)} + \mathrm{j} \omega_{\mathrm{\varphi_s}} \cdot \underline{\psi}_{\mathrm{s}}(t).
    \end{equation}

    Due to the steady state, the change of the flux amplitude $\left(\frac{\mathrm{d}}{\mathrm{d}t}\psi(t)\right)$ is zero. The flux vector rotates with the speed $\omega$, therefore, the flux vector is determined as follows:
    \begin{equation}
        \underline{\psi}_{\mathrm{s}}(t)
        = \frac{\SI{283}{\volt} \cdot e^{\mathrm{j}44.4^{\circ}}}{\SI{314.16}{\per\second}\cdot e^{\mathrm{j}90^{\circ}}}
        = \SI{0.90}{\volt\second} \cdot e^{-\mathrm{j}\SI{45.6}{\degree}}.
    \end{equation}


    \begin{solutionfigure}
        \centering
        \includegraphics{ex05/vectorDiagram_IM.pdf}
        \caption{Resulting vector diagram of the induction machine in steady state. The scala for the voltage is 1 cm $\widehat{=}~\SI{100}{\volt} $, for the flux linkage 1 cm $\widehat{=}~\SI{1}{\volt\second}$ and for the current 1 cm $\widehat{=}~\SI{10}{\ampere}$.}
        \label{fig:vectorDiagram_IM}
    \end{solutionfigure}


\end{solutionblock}